\subsection{Color Channel Correlation}
Trong ảnh tự nhiên, ba kênh màu R, G, B có tương quan do ánh sáng, bề mặt vật thể, cảm biến và pipeline xử lý ảnh. Ảnh AI có thể sinh kênh màu theo cách không hoàn toàn giống camera thật, dẫn đến tương quan bất thường hoặc quá đều.

\begin{longtable}{p{0.06\textwidth}p{0.22\textwidth}p{0.30\textwidth}p{0.24\textwidth}p{0.15\textwidth}}
\caption{Tóm tắt 5 nhóm tín hiệu hiệu quả của SourceVerify}\label{tab:signal-groups}\\
\toprule
\textbf{STT} & \textbf{Nhóm tín hiệu} & \textbf{Method đại diện} & \textbf{Dấu hiệu phát hiện} & \textbf{Hạn chế chính} \\
\midrule
\endfirsthead
\toprule
\textbf{STT} & \textbf{Nhóm tín hiệu} & \textbf{Method đại diện} & \textbf{Dấu hiệu phát hiện} & \textbf{Hạn chế chính} \\
\midrule
\endhead
1 & Nguồn gốc tệp & Metadata Analysis & Phần mềm tạo ảnh, EXIF, kích thước, tên tệp & Dễ bị xoá/sửa \\
2 & Miền tần số & Spectral Nyquist, Multi-scale Reconstruction & Đỉnh tần số, artifact upsampling, sai số tái tạo bất thường & Bị ảnh hưởng bởi resize \\
3 & Nhiễu/kết cấu & Noise Residual, Gradient Micro-Texture & Nhiễu quá sạch, thiếu micro-texture tự nhiên & Nhạy với nén/filter \\
4 & Biên cạnh/thống kê & Edge Coherence, Benford's Law & Biên quá mượt, phân bố gradient lệch tự nhiên & Không kết luận đơn lẻ \\
5 & Màu sắc/nén ảnh & Chromatic Aberration, DCT Block Artifacts, Color Channel Correlation & Sai lệch màu, vết nén JPEG, tương quan RGB bất thường & Phụ thuộc chỉnh sửa ảnh \\
\bottomrule
\end{longtable}
